% ============================================================
%  A Mathematical Theory of Market World Models
%  HongJin HE, Alpha Flow Research
%  Target: Mathematical Finance / SIAM J. Financial Mathematics
% ============================================================

\documentclass[12pt,a4paper]{article}

% ----- Packages -----
\usepackage{amsmath, amsthm, amssymb, amsfonts}
\usepackage{mathtools}
\usepackage[margin=1in]{geometry}
\usepackage{hyperref}
\usepackage{natbib}
\usepackage{bm}
\usepackage{bbm}
\usepackage{enumitem}
\usepackage{booktabs}
\usepackage{array}
\usepackage{xcolor}
\usepackage{setspace}
\onehalfspacing

% ----- Theorem Environments -----
\theoremstyle{plain}
\newtheorem{theorem}{Theorem}[section]
\newtheorem{proposition}[theorem]{Proposition}
\newtheorem{lemma}[theorem]{Lemma}
\newtheorem{corollary}[theorem]{Corollary}

\theoremstyle{definition}
\newtheorem{definition}[theorem]{Definition}
\newtheorem{assumption}[theorem]{Assumption}
\newtheorem{example}[theorem]{Example}
\newtheorem{remark}[theorem]{Remark}

% ----- Macros -----
\newcommand{\R}{\mathbb{R}}
\newcommand{\E}{\mathbb{E}}
\newcommand{\PP}{\mathbb{P}}
\newcommand{\F}{\mathcal{F}}
\newcommand{\Hil}{\mathcal{H}}
\newcommand{\Z}{\mathcal{Z}}
\newcommand{\A}{\mathcal{A}}
\newcommand{\Prob}{\mathcal{P}}
\newcommand{\Law}{\mathrm{Law}}
\newcommand{\tr}{\mathrm{tr}}
\newcommand{\diag}{\mathrm{diag}}
\newcommand{\Var}{\mathrm{Var}}
\newcommand{\Cov}{\mathrm{Cov}}
\newcommand{\1}{\mathbbm{1}}
\newcommand{\eps}{\varepsilon}
\newcommand{\vp}{\varphi}
\newcommand{\norm}[1]{\left\|#1\right\|}
\newcommand{\abs}[1]{\left|#1\right|}
\newcommand{\inner}[2]{\left\langle #1,\, #2 \right\rangle}
\newcommand{\WD}[2]{W_2(#1,#2)}  % Wasserstein-2 distance
\newcommand{\KL}[2]{D_{\mathrm{KL}}(#1 \| #2)}
\DeclareMathOperator*{\argmin}{arg\,min}
\DeclareMathOperator*{\argmax}{arg\,max}
\DeclareMathOperator{\Bias}{Bias}
\DeclareMathOperator{\src}{src}
\DeclareMathOperator{\tgt}{tgt}

% ============================================================
\begin{document}
% ============================================================

\title{\Large\bfseries A Mathematical Theory of Market World Models:\\[6pt]
Hierarchical Mean-Field Dynamics, Dual Stochastic Decomposition,\\
and Financial Event Operator Algebras}

\author{
  HongJin HE\\[6pt]
  \small The Hong Kong University of Science and Technology \quad
         Stanford University (IHP)\\[3pt]
  \small \texttt{hhjhk@stanford.edu}
}

\date{\small July 2026 \quad \textit{(Preprint — submitted to arXiv:q-fin.MF)}}

\maketitle

% ============================================================
\begin{abstract}
We develop a rigorous mathematical theory for \emph{market world models}—
probabilistic dynamical systems that learn and simulate the full evolutionary
law of financial markets rather than extracting statistical correlations through
factor models.  Our framework rests on four interlocking contributions.

\textbf{(i) Dual Stochastic Decomposition.} We prove that market noise admits
a canonical Lévy–Itô decomposition into two orthogonal components: a
\emph{physical noise} driven by a Brownian motion (capturing fundamental
information uncertainty) and a \emph{behavioral noise} driven by a pure-jump
Lévy process (capturing sentiment and herding effects).  We derive a
Cramér–Rao-type lower bound showing that prediction variance is bounded below
by a sum of two irreducible terms arising from each component—a market-structure
analogue of the Heisenberg uncertainty principle.

\textbf{(ii) Financial Event Operator Algebra.} Every corporate or macroeconomic
event (IPO, merger, interest rate shock, etc.) is modeled as an affine
operator on the market state space.  We classify all such operators by their
algebraic mode and prove that the full event algebra carries the structure of a
\emph{topological groupoid}—a generalisation of a semigroup required because
mergers and spin-offs alter the dimension of the state space.

\textbf{(iii) Hierarchical Mean-Field Theory.} Real financial markets exhibit
simultaneous mean-field game (MFG) and mean-field control (MFC) structures at
different scales: nations interact as MFG agents while firms within each nation
face a mixed MFG/MFC problem regulated by a central authority.  We formulate a
two-level McKean–Vlasov system and prove existence and uniqueness of a
\emph{hierarchical Nash equilibrium} under Lipschitz, boundedness, and
Lasry–Lions monotonicity conditions.

\textbf{(iv) Stochastic Stability.} Using the stochastic Lyapunov method for
systems driven by the combined Brownian-Lévy noise, we prove that the market
world model admits a unique invariant measure and establish sharp exponential
convergence rates. Market bubbles and crises correspond precisely to local
violations of the Lyapunov condition.

Taken together, these results constitute, to the best of our knowledge,
the \emph{first} systematic mathematical theory for market world models:
a self-contained, axiomatically grounded framework that connects the
world-model paradigm of machine learning to the rigorous traditions of
stochastic analysis and mean-field game theory in quantitative finance.

\smallskip
\noindent\textbf{Keywords:} world model, mean-field game, stochastic
differential equation, Lévy process, groupoid, Lyapunov stability,
market microstructure, prediction lower bound.

\smallskip
\noindent\textbf{MSC 2020:} 91G99, 60G51, 49N80, 22A22, 60H10.
\end{abstract}

\tableofcontents

% ============================================================
\section{Introduction}
% ============================================================

\subsection{Motivation}

The dominant paradigm in quantitative finance rests on two pillars: factor
models, which represent asset returns as linear combinations of explanatory
signals, and strategy models, which exploit persistent anomalies for
profit~\citep{Fama1970,Ross1976}. Despite their practical success, these
approaches share a fundamental limitation: they learn \emph{correlations} in
historical data without learning the underlying \emph{causal dynamics} of
financial markets.

A market \emph{world model} takes a different approach. Inspired by
the world-model framework of \citet{HaSchmidhuber2018} in reinforcement
learning—where an agent maintains an internal model of its environment's
temporal dynamics—a market world model aims to learn the full evolutionary law
governing market states, enabling counterfactual reasoning and scenario
propagation unavailable to correlation-based methods.

Despite growing empirical interest~\citep{AwmlgLiverpool2026}, \emph{no
systematic mathematical theory} has been established for market world models.
Existing work remains either purely algorithmic (lacking theoretical
guarantees) or limited to single-scale models (ignoring the hierarchical
multi-agent structure fundamental to financial markets).
This paper provides the missing foundation.

\subsection{Contributions}

Our four main contributions, each embodied in one or more formal theorems, are:

\begin{enumerate}[label=(\roman*)]
\item \textbf{Dual Stochastic Decomposition} (Section~\ref{sec:noise}).
  Theorem~\ref{thm:noise_decomp} establishes the canonical Lévy--Itô
  decomposition of market noise into physical and behavioral components.
  Theorem~\ref{thm:cramer_rao} provides a prediction lower bound—an
  information-theoretic impossibility result showing that no estimator can
  eliminate both components of forecast uncertainty.

\item \textbf{Financial Event Operator Algebra} (Section~\ref{sec:events}).
  Definition~\ref{def:operator} introduces affine event operators.
  Theorem~\ref{thm:groupoid} proves that the full algebra of financial events
  constitutes a topological groupoid.

\item \textbf{Hierarchical Mean-Field Equilibrium} (Section~\ref{sec:mfg}).
  Theorem~\ref{thm:hier_mfg} is the main result of the paper: existence and
  uniqueness of a hierarchical Nash equilibrium for the two-level McKean--Vlasov
  system with a Lasry--Lions monotonicity structure.

\item \textbf{Stochastic Lyapunov Stability} (Section~\ref{sec:stability}).
  Theorem~\ref{thm:lyapunov} establishes existence of a unique invariant
  measure and exponential convergence; Proposition~\ref{prop:bubble} identifies
  the analytic signature of market bubbles and crashes.
\end{enumerate}

\subsection{Related Work}

\paragraph{Mean-field games in finance.}
The MFG framework of \citet{LasryLions2007} and \citet{HuangMalhame2006}
has been applied to price impact~\citep{CarmonaLacker2015},
systemic risk~\citep{CarmonaMouzouni2023}, and optimal
liquidation~\citep{CarmonaDelarue2018book}. Multi-level MFG has been studied
in the context of inter-bank lending~\citep{CarmonaSun2018} but without the
hierarchical coupling structure we introduce here.

\paragraph{Originality statement.}
To the best of our knowledge, this paper is the \emph{first} to:
(a)~formally define a market world model as a probability-measure-valued
map on an infinite-dimensional state space;
(b)~prove a Cramér--Rao-type lower bound that identifies two structurally
distinct and irreducible components of market prediction uncertainty;
(c)~model the full algebra of financial events as a topological groupoid,
capturing the dimension-changing nature of mergers and spin-offs;
(d)~establish existence and uniqueness of a hierarchical Nash equilibrium
in a two-level McKean--Vlasov system that simultaneously exhibits
MFG and MFC structures at different scales.
Prior work on world models in finance is either empirical without theory
(\citealt{AwmlgLiverpool2026}) or models only single-scale dynamics
(\citealt{CarmonaDelarue2018book}).

\paragraph{Jump processes in finance.}
Lévy-driven asset prices are classical~\citep{ContTankov2004,Kou2002}.
The decomposition of price innovations into a diffusion and a jump component
has been studied for estimation purposes~\citep{AitSahaliaTodorov2009}, but
not as a \emph{structural} duality between physical and behavioral uncertainty.

\paragraph{World models and reinforcement learning.}
Since \citet{HaSchmidhuber2018}, world models have been central to
model-based RL~\citep{Hafner2020Dreamer,Hafner2023DreamerV3}. Their application
to finance is nascent: \citet{AwmlgLiverpool2026} show that a generative world
model improves portfolio optimization; \citet{MorganJPMaGAN} apply GANs to
order-book simulation. No prior work provides a unified mathematical theory.

\paragraph{Groupoids in physics and mathematics.}
Groupoids generalize groups by allowing morphisms between different objects.
Their use in noncommutative geometry~\citep{Connes1994} and geometric
mechanics~\citep{Weinstein1996} is well-established; their appearance in
financial modeling is novel.

% ============================================================
\section{Probabilistic Foundations}
\label{sec:prelim}
% ============================================================

\subsection{Probability Space and Filtration}

Throughout the paper, we fix a complete probability space
$(\Omega, \F, \PP)$ equipped with a right-continuous, complete filtration
$\mathbb{F} = (\F_t)_{t \in [0,T]}$ for a fixed horizon $T \in (0,\infty)$.

We work with two independent sources of randomness:
\begin{itemize}
  \item A standard $\R^m$-valued Brownian motion $W = (W_t)_{t \geq 0}$, and
  \item A Poisson random measure $N^\eta$ on $[0,T] \times \R^k$ with
        (non-zero) $\sigma$-finite intensity measure $\lambda \otimes \nu^\eta$,
        where $\lambda$ denotes Lebesgue measure. The compensated measure is
        $\tilde{N}^\eta(dt,dz) := N^\eta(dt,dz) - \nu^\eta(dz)\,dt$.
\end{itemize}
Both $W$ and $N^\eta$ are adapted to $\mathbb{F}$ and mutually independent.

\subsection{Itô Stochastic Integral: The Correct Formulation}

A critical structural observation distinguishes our framework from ordinary
differential equation (ODE) based approaches:

\begin{proposition}[Non-differentiability of Brownian paths]
\label{prop:nondiff}
With probability one, $t \mapsto W_t$ is nowhere differentiable. Consequently,
the expression $\sigma_t \cdot (dW_t/dt)$ is meaningless pointwise, and the
shorthand notation $dS_t = \mu_t\,dt + \sigma_t\,dW_t$ represents an
\emph{integral equation}, not a differential equation.
\end{proposition}

\begin{proof}
For any $t$, the increment $h^{-1}(W_{t+h} - W_t) \sim \mathcal{N}(0,h^{-1})$.
As $h \to 0$, the variance $h^{-1} \to \infty$, so no $L^2$ limit exists.
By the Paley--Wiener--Zygmund theorem, the set of $\omega$ for which $W_\cdot(\omega)$
is differentiable at any single point has measure zero; a countable union
argument extends this to all of $[0,T]$.
\end{proof}

The precise meaning of the market evolution equation is therefore:

\begin{definition}[Itô integral form]
\label{def:ito_integral}
The market state process satisfies the \emph{stochastic integral equation}:
\begin{equation}
\label{eq:full_sde}
\mathbf{S}_t = \mathbf{S}_0
  + \int_0^t \mu(\mathbf{S}_u, u)\,du
  + \int_0^t \sigma_\tau(\mathbf{S}_u)\,dW_u
  + \int_0^t\!\!\int_{\R^k} z\,\tilde{N}^\eta(du,dz)
  + \sum_{\tau_w \leq t} (T_w - I)\mathbf{S}_{\tau_w^-},
\end{equation}
where the second and third Itô integrals are defined as $L^2(\Omega)$ limits of
simple predictable approximants, and the last sum runs over all financial event
times $\tau_w \in [0,t]$.
\end{definition}

The following isometry is fundamental:

\begin{proposition}[Itô isometry]
For any $\mathbb{F}$-predictable process $\sigma_\cdot$ with
$\E[\int_0^T \|\sigma_t\|^2\,dt] < \infty$,
\[
\E\!\left[\left\|\int_0^t \sigma_u\,dW_u\right\|^2\right]
= \E\!\left[\int_0^t \|\sigma_u\|^2\,du\right].
\]
The adaptedness requirement $\sigma_u \in \F_u$ encodes the economic
constraint of \emph{no look-ahead}: trading signals may not depend on
future information.
\end{proposition}

\subsection{Lévy Processes and the Lévy--Khintchine Formula}

\begin{definition}[Lévy process]
A process $L = (L_t)_{t \geq 0}$ is a \emph{Lévy process} if
(i)~$L_0 = 0$ a.s., (ii)~$L$ has independent, stationary increments,
and (iii)~$L$ is continuous in probability.
\end{definition}

\begin{theorem}[Lévy--Khintchine, \citealt{Sato1999}]
\label{thm:LK}
The characteristic function of $L_t$ has the form
$\E[e^{i\xi^\top L_t}] = e^{-t\Psi(\xi)}$, where the Lévy exponent is:
\[
\Psi(\xi)
= -i b^\top \xi
  + \tfrac{1}{2}\xi^\top C\xi
  + \int_{\R^k}\!\left(1 - e^{i\xi^\top z} + i\xi^\top z\,\1_{|z|\leq 1}\right)
    \nu(dz),
\]
with \emph{drift} $b \in \R^k$, \emph{diffusion matrix}
$C \in \R^{k\times k}$ positive semi-definite, and
\emph{Lévy measure} $\nu$ satisfying
$\int_{\R^k}(1 \wedge |z|^2)\,\nu(dz) < \infty$.
The triple $(b,C,\nu)$ is the \emph{Lévy characteristic} and determines $L$
in law.
\end{theorem}

\subsection{Market State Space}

\begin{definition}[Single-asset state vector]
\label{def:state_vector}
The state of a single asset at time $t$ is the vector
$s_t = (p_t, v_t, \ell_t, \kappa_t, \iota_t)^\top \in \R^5$, where:
\begin{center}
\begin{tabular}{clll}
\toprule
Coord. & Symbol & Economic meaning & Domain \\
\midrule
1 & $p_t$ & Log-price & $\R$ \\
2 & $v_t$ & Log-trading volume & $\R_+$ \\
3 & $\ell_t$ & Leverage ratio & $[0,\infty)$ \\
4 & $\kappa_t$ & Outstanding shares (log) & $\R_+$ \\
5 & $\iota_t$ & Public information disclosure level & $[0,1]$ \\
\bottomrule
\end{tabular}
\end{center}
\end{definition}

For $n$ assets, the joint state is $\mathbf{S}_t \in \R^{5n}$.
In the continuum limit ($n \to \infty$) relevant for theoretical analysis, we
work in the separable Hilbert space
$\Hil = L^2([0,1]\times[0,T],\R^5)$ with inner product
$\inner{f}{g} = \int_0^1\!\int_0^T f(u,t)\cdot g(u,t)\,dt\,du$.

% ============================================================
\section{The World Model as an Optimal Estimation Problem}
\label{sec:world_model}
% ============================================================

\begin{definition}[Market world model]
\label{def:world_model}
A \emph{market world model} with parameters $\theta$ is a measurable map
\[
\mathcal{M}_\theta : \F_t \times (0,\infty) \to \Prob(\Hil)
\]
such that $\mathcal{M}_\theta(\F_t, h) \approx \PP(\mathbf{S}_{t+h} \in \cdot \mid \F_t)$.
\end{definition}

\begin{theorem}[Prediction error decomposition]
\label{thm:pred_decomp}
For any $\F_t$-measurable estimator $\hat{\mathbf{S}}_{t+h}$ of $\mathbf{S}_{t+h}$,
\begin{equation}
\label{eq:pred_decomp}
\E\!\left[\norm{\hat{\mathbf{S}}_{t+h} - \mathbf{S}_{t+h}}_\Hil^2\right]
= \norm{\Bias_\theta}^2
  + \underbrace{\Var_\tau(h)}_{\text{physical}}
  + \underbrace{\Var_\eta(h)}_{\text{behavioral}},
\end{equation}
where the last two terms are strictly positive and model-independent
(see Theorem~\ref{thm:cramer_rao}).
\end{theorem}

\begin{proof}
The bias--variance decomposition is standard. For the physical variance, write
$\mathbf{S}_{t+h} = \hat{\mathbf{S}}_{t+h} + (\mathbf{S}_{t+h} - \E[\mathbf{S}_{t+h}|\F_t])
+ (\E[\mathbf{S}_{t+h}|\F_t] - \hat{\mathbf{S}}_{t+h})$ and expand using
$\E[M_t(A_t - B)] = 0$ for a martingale $M$ and $\F_t$-measurable $B$.
The decomposition into physical and behavioral components follows from the
independence established in Theorem~\ref{thm:noise_decomp}.
\end{proof}

% ============================================================
\section{Dual Stochastic Decomposition}
\label{sec:noise}
% ============================================================

\subsection{Two Sources of Market Uncertainty}

Standard models (e.g., Black--Scholes) conflate two qualitatively different
sources of market randomness into a single Brownian motion.
We argue this conflation is structurally incorrect:

\begin{itemize}
\item \textbf{Physical noise} $\tau_t$: Arises from fundamental information
  asymmetry, parameter estimation error, and market microstructure frictions.
  It accumulates \emph{continuously} along trading time—analogous to thermal
  noise in physical systems.

\item \textbf{Behavioral noise} $\eta_t$: Arises from investor sentiment,
  herding, panic selling, and euphoric buying. It manifests as
  \emph{discontinuous jumps}—sudden regime shifts that cannot be generated by
  any continuous Brownian path.
\end{itemize}

\begin{assumption}[Dual noise structure]
\label{ass:dual_noise}
The noise in equation~\eqref{eq:full_sde} decomposes as:
\begin{equation}
\label{eq:dual_noise}
\text{Noise}_t = \underbrace{\sigma_\tau(\mathbf{S}_t)\,dW_t^{(\tau)}}_{\text{physical}} + \underbrace{dJ_t^{(\eta)}}_{\text{behavioral}},
\end{equation}
where $W^{(\tau)}$ is a standard Brownian motion and
$J^{(\eta)}_t = \int_0^t\!\int_{\R^k} z\,\tilde{N}^\eta(ds,dz)$ is an
independent pure-jump Lévy process with Lévy measure $\nu^\eta$ satisfying
$\int_{|z|>1}|z|^2\,\nu^\eta(dz) < \infty$ and
$[W^{(\tau)}, J^{(\eta)}]_t = 0$ a.s.
\end{assumption}

\subsection{Canonical Decomposition Theorem}

\begin{theorem}[Dual Lévy--Itô decomposition]
\label{thm:noise_decomp}
Under Assumption~\ref{ass:dual_noise}, the total noise process $L_t :=
\sigma_\tau W_t^{(\tau)} + J_t^{(\eta)}$ has a unique Lévy characteristic
$(b, C, \nu)$ that decomposes as:
\begin{equation}
\label{eq:levy_char_decomp}
b = b_\tau + b_\eta, \quad
C = \Sigma_\tau := \sigma_\tau\sigma_\tau^\top, \quad
\nu = \nu^\eta,
\end{equation}
with $\Sigma_\tau$ (the diffusion matrix) and $\nu^\eta$ (the jump intensity
measure) \emph{uniquely identified} by the characteristic function:
\begin{equation}
\label{eq:char_fn}
\E\!\left[e^{i\xi^\top L_t}\right]
= \exp\!\left(-t\!\left[
    \tfrac{1}{2}\xi^\top\Sigma_\tau\xi
    + \int_{\R^k}\!\left(1 - e^{i\xi^\top z} + i\xi^\top z\,\1_{|z|\leq1}\right)
      \nu^\eta(dz)
  \right]\right).
\end{equation}
In particular, $\Sigma_\tau \neq 0$ and $\nu^\eta \neq 0$ are \emph{orthogonal}
in the sense that no jump measure can generate a continuous Gaussian component.
\end{theorem}

\begin{proof}
By independence of $W^{(\tau)}$ and $J^{(\eta)}$, the characteristic function
of $L_t$ factors as
\[
\E[e^{i\xi^\top L_t}]
= \E[e^{i\xi^\top \sigma_\tau W_t}]\cdot\E[e^{i\xi^\top J_t^\eta}]
= e^{-t\Psi_\tau(\xi)}\cdot e^{-t\Psi_\eta(\xi)},
\]
where $\Psi_\tau(\xi) = \frac{1}{2}\xi^\top\Sigma_\tau\xi$ (Gaussian) and
$\Psi_\eta(\xi) = \int(1-e^{i\xi^\top z}+i\xi^\top z\,\1_{|z|\leq 1})\nu^\eta(dz)$
(pure-jump Lévy). Adding the two Lévy exponents gives~\eqref{eq:char_fn}.
Uniqueness of the Lévy characteristic (Theorem~\ref{thm:LK}) and the
canonical decomposition of Lévy processes~\citep{Sato1999} establish that
$\Sigma_\tau$ and $\nu^\eta$ are uniquely determined. The orthogonality
statement reflects the Lévy--Khintchine classification: the Gaussian part is
determined by $C = \Sigma_\tau$ alone, the jump part by $\nu^\eta$ alone;
no assignment of mass to $\nu^\eta$ at any point in $\R^k$ can contribute
to $C$.
\end{proof}

\begin{remark}[Heavy tails of behavioral noise]
Empirical evidence~\citep{ContBouchaud2000,LuxMarchesi1999} supports
$\nu^\eta(|z|>x) \sim x^{-\alpha}$ for $\alpha \in (1,2)$ (heavy-tailed
but with finite variance), consistent with the power-law tails observed in
financial return distributions and inconsistent with any Gaussian model.
\end{remark}

\subsection{Prediction Lower Bound}

\begin{theorem}[Cramér--Rao-type prediction lower bound]
\label{thm:cramer_rao}
Let $\mathbf{S}$ satisfy~\eqref{eq:full_sde} under Assumption~\ref{ass:dual_noise}.
For any $\F_t$-measurable unbiased estimator $\hat{\mathbf{S}}_{t+h}$ of
$\mathbf{S}_{t+h}$ and any horizon $h > 0$:
\begin{equation}
\label{eq:CR_bound}
\Var\!\left(\hat{\mathbf{S}}_{t+h} - \mathbf{S}_{t+h}\right)
\geq \underbrace{\sigma_\tau^2\,h}_{\text{physical lower bound}}
   + \underbrace{\lambda_\eta\,m_2^\eta\,h}_{\text{behavioral lower bound}}
\eqqcolon \mathrm{CR}(h),
\end{equation}
where $\sigma_\tau^2 = \tr(\Sigma_\tau)$, $\lambda_\eta > 0$ is the jump
intensity, and $m_2^\eta = \int_{\R^k}|z|^2\,\nu^\eta(dz)$.
\end{theorem}

\begin{proof}
\textbf{Step 1 (Physical component).}
Conditional on $\F_t$, the increment $\mathbf{S}_{t+h}^{(\tau)} - \mathbf{S}_t^{(\tau)}
:= \int_t^{t+h}\sigma_\tau\,dW_u$ is Gaussian with mean zero and covariance
$\Sigma_\tau h$.  For any unbiased estimator $\hat{A}$ of a Gaussian parameter
with Fisher information $I$, the standard Cramér--Rao inequality gives
$\Var(\hat{A}) \geq I^{-1}$.  Here $I^{-1} = \Sigma_\tau h$, so
$\Var(\hat{\mathbf{S}}^{(\tau)}_{t+h}) \geq \tr(\Sigma_\tau)h = \sigma_\tau^2 h$.

\textbf{Step 2 (Behavioral component).}
Let $K_{t,t+h} = N^\eta([t,t+h]\times\R^k)$ be the number of jumps in
$[t,t+h]$. Given $K_{t,t+h} = n$, the jump sizes are i.i.d.\ with common
distribution $\nu^\eta(dz)/\lambda_\eta$, so the total jump contribution
has variance $n \cdot m_2^\eta / \lambda_\eta$.  Integrating over
$K_{t,t+h} \sim \mathrm{Poisson}(\lambda_\eta h)$:
\[
\Var\!\left(\int_t^{t+h}\!\!\int z\,\tilde{N}^\eta(du,dz)\right)
= \lambda_\eta h \cdot m_2^\eta.
\]
Since $\hat{\mathbf{S}}_{t+h}$ cannot predict the future jump realisation
without observing $N^\eta$ beyond time $t$, this variance is irreducible.

\textbf{Step 3 (Independence).}
By Assumption~\ref{ass:dual_noise}, physical and behavioral increments are
independent, so their variances add:
$\Var(\mathbf{S}_{t+h}^{(\tau)} + J_{t+h}^{(\eta)} - J_t^{(\eta)})
= \sigma_\tau^2 h + \lambda_\eta m_2^\eta h = \mathrm{CR}(h)$.
\end{proof}

\begin{remark}[Analogy with quantum mechanics]
Theorem~\ref{thm:cramer_rao} is the financial analogue of the Heisenberg
uncertainty principle: even with perfect knowledge of $\F_t$ and an optimal
model, the sum of prediction variances across physical and behavioral channels
cannot fall below $\mathrm{CR}(h)$.  This bound is a structural feature of
the market, not a limitation of methodology.
\end{remark}

\begin{corollary}[Time-varying temperature parameter]
The scalar temperature parameter $\tau_t$ of the world model~\citep{HaSchmidhuber2018}
has the closed-form expression:
\[
\tau_t = \sqrt{\sigma_\tau^2(\mathbf{S}_t) + \lambda_\eta(t)\,m_2^\eta(t)},
\]
and is a $\mathbb{F}$-adapted, time-varying quantity determined by local
noise intensities.  Training the world model to match the observed local
variance provides a principled, data-driven estimate of $\tau_t$.
\end{corollary}

% ============================================================
\section{Financial Event Operator Algebra}
\label{sec:events}
% ============================================================

\subsection{Affine Event Operators}

\begin{definition}[Affine event operator]
\label{def:operator}
For each financial event type $w \in \mathcal{W}$, the associated
\emph{event operator} is the affine map:
\begin{equation}
T_w(s) = A_w\,s + b_w + \Sigma_w\,\varepsilon_w,
\quad \varepsilon_w \sim \mathcal{N}(0, I),
\end{equation}
where $A_w \in \R^{d'\times d}$ (with $d' = d$ or $d'\neq d$ depending on
the event type), $b_w \in \R^{d'}$, and
$\Sigma_w \in \R^{d'\times d'}$ encodes event-specific announcement uncertainty.
\end{definition}

\subsection{Classification of Financial Events}

Financial events fall into three algebraic modes based on how they act on the
state space:

\begin{definition}[Three modes of action]
\label{def:modes}
\begin{enumerate}[label=(\Roman*)]
\item \textbf{Local endomorphism} ($d' = d$, acts on one asset):
  stock split, share repurchase, earnings announcement, dividend payment.
  $A_w \in \R^{d\times d}$; the state space dimension is preserved.
\item \textbf{Global tensor action} (acts on all $n$ assets via Kronecker product):
  central bank rate change, systemic shock, sector-wide regulation.
  $T_w^{\mathrm{global}} = \Lambda_w \otimes I_d$, with
  $\Lambda_w \in \R^{n\times n}$ encoding heterogeneous sensitivities.
\item \textbf{Pairwise morphism} ($d'\neq d$, changes entity count):
  merger \& acquisition ($n\to n-1$), spin-off / IPO ($n\to n+1$).
  $A_w$ is \emph{non-square}; the state space dimension changes.
\end{enumerate}
\end{definition}

\begin{proposition}[Information irreversibility]
\label{prop:info_irrev}
For any Type-I event $w$ with natural inverse $w^{-1}$ (e.g., stock split
followed by reverse split),
\[
T_{w^{-1}} \circ T_w = I + \mathcal{E}^\mathrm{info}_{w},
\]
where $\mathcal{E}^\mathrm{info}_{w} \neq 0$ is an operator acting non-trivially
on the information subspace $\{\iota\text{-coordinate}\}$ of the state vector.
Even when $A_w$ is invertible, the event releases information that cannot be
unannounced: matrix invertibility does not imply physical reversibility.
\end{proposition}

\begin{proof}
Let $s = (p, v, \ell, \kappa, \iota)^\top$. For a stock split with ratio
$r > 1$: $A_w = \diag(1,1,1,r,1)$, $b_w = 0$. The split also increases the
disclosure level: $\iota \mapsto \min(\iota + \delta_w, 1)$ for some
$\delta_w > 0$. Hence $b_w$ has a positive entry in the $\iota$-coordinate.
The inverse operation $A_{w^{-1}} = \diag(1,1,1,r^{-1},1)$ restores
$\kappa$ but leaves $\iota$ elevated, so
$(T_{w^{-1}} \circ T_w)(s)_\iota = \iota + \delta_w \neq \iota$.
\end{proof}

\subsection{Groupoid Structure of the Event Algebra}

Type-III events (mergers and spin-offs) change the state space dimension $n$,
making the operator $A_w$ non-square. This prevents the collection of all
event operators from forming a semigroup (which requires a single carrier set).
The correct algebraic structure is a \emph{groupoid}.

\begin{definition}[Financial event groupoid]
\label{def:groupoid}
Define the \emph{financial event groupoid}
$\mathcal{G}_{\mathrm{fin}} = (\mathcal{W}, \circ, \src, \tgt, \mathrm{id})$ by:
\begin{itemize}
  \item \textbf{Objects}: $\mathrm{Ob}(\mathcal{G}) = \{\Hil_S : S \text{ is a
        valid $n$-asset market configuration, } n \in \N\}$.
  \item \textbf{Morphisms}: For $T_w \in \mathrm{Hom}(\Hil_S, \Hil_{S'})$,
        $\src(T_w) = \Hil_S$ and $\tgt(T_w) = \Hil_{S'}$.
  \item \textbf{Composition}: $T_{w_2} \circ T_{w_1}$ is defined if and only if
        $\src(T_{w_2}) = \tgt(T_{w_1})$.
  \item \textbf{Identities}: $\mathrm{id}_{\Hil_S} = I_{d\cdot|S|}$.
  \item \textbf{Topology}: Endow each $\mathrm{Hom}(\Hil_S, \Hil_{S'})$ with
        the operator norm topology; the whole space with the disjoint-union
        topology.
\end{itemize}
\end{definition}

\begin{theorem}[Groupoid structure]
\label{thm:groupoid}
$\mathcal{G}_{\mathrm{fin}}$ is a topological groupoid.
\end{theorem}

\begin{proof}
We verify the four groupoid axioms:

\textit{(i) Well-defined composition.} Given $T_{w_1}: \Hil_S \to \Hil_{S'}$
and $T_{w_2}: \Hil_{S'} \to \Hil_{S''}$, their composition
$T_{w_2} \circ T_{w_1}: \Hil_S \to \Hil_{S''}$ is the standard
operator composition, which is well-defined in Banach spaces whenever the
intermediate spaces match. The dimensionality compatibility condition
$\src(T_{w_2}) = \tgt(T_{w_1})$ is exactly this matching requirement.

\textit{(ii) Associativity.} Operator composition in Banach spaces is
associative: $(T_{w_3}\circ T_{w_2})\circ T_{w_1}
= T_{w_3}\circ(T_{w_2}\circ T_{w_1})$ whenever both sides are defined.

\textit{(iii) Identity morphisms.} For each $\Hil_S$, the identity
operator $\mathrm{id}_{\Hil_S} = I_{d|S|}$ satisfies
$T_w \circ \mathrm{id}_{\Hil_S} = T_w$ and
$\mathrm{id}_{\Hil_{S'}} \circ T_w = T_w$ for any
$T_w \in \mathrm{Hom}(\Hil_S, \Hil_{S'})$.

\textit{(iv) Inverses (partial).} For Type-I and Type-II events with
invertible $A_w$ (i.e., $\det(A_w) \neq 0$), the inverse $T_w^{-1}$ exists
as an operator and $T_w \circ T_w^{-1} = \mathrm{id}_{\Hil_S}$
(up to the information residual in Proposition~\ref{prop:info_irrev}).
For Type-III events (non-square $A_w$), the Moore--Penrose pseudoinverse
$T_w^+$ satisfies $T_w T_w^+ T_w = T_w$ but is not a full inverse—this
is the natural partial-inverse structure of a groupoid.

\textit{(v) Continuity.} Under the operator norm topology,
$(T_{w_1}, T_{w_2}) \mapsto T_{w_2} \circ T_{w_1}$ is continuous
(composition of bounded operators is jointly continuous in the norm topology).
Similarly, $T_w \mapsto T_w^+$ is continuous on the set of operators
with fixed rank.  Hence $\mathcal{G}_{\mathrm{fin}}$ is a topological groupoid.
\end{proof}

\begin{corollary}[Endomorphism monoid as a sub-structure]
The collection of Type-I events $\mathcal{W}_I \subset \mathcal{W}$ acting
on a fixed $n$-asset space forms an \emph{endomorphism monoid} with
identity $I_{5n}$. For Type-I events with $\det(A_w) \neq 0$ this monoid
is a group.
\end{corollary}

% ============================================================
\section{The E-Game-C Architecture}
\label{sec:egc}
% ============================================================

\subsection{Limitations of the Original VMC World Model}

The V--M--C architecture of \citet{HaSchmidhuber2018} consists of:
a Variational AutoEncoder (\textbf{V}) that compresses high-dimensional
observations to a latent code $z_t$;
a recurrent world model (\textbf{M}) that predicts $z_{t+1}$ from $(z_t, a_t)$;
and a linear controller (\textbf{C}) trained by evolution strategies.

Direct application to finance faces three structural obstacles:
(a) the input is heterogeneous (prices, news, event sequences), not pixels;
(b) the ``environment'' is not fixed—it is the aggregate outcome of all
agents' decisions, creating a reflexive feedback loop;
(c) the world model module $M$ implicitly assumes the environment is a
fixed dynamical system, ignoring the multi-agent strategic structure
that is fundamental to financial markets.

\subsection{Encoder Module E}

\begin{definition}[Information encoder]
\label{def:encoder}
The encoder $\mathcal{E}_\phi: \F_t \to \Z$ maps the information set at time
$t$ to a latent code $z_t \in \Z \subset \R^{d_z}$ via a
variational autoencoder~\citep{KingmaWelling2014}:
\[
z_t = \mathcal{E}_\phi(\mathbf{x}_t), \quad
\mathbf{x}_t = (p_t, v_t, \text{news}_t, \text{events}_t, \ldots) \in \F_t,
\]
trained by maximising the ELBO:
\[
\mathcal{L}_{\mathrm{ELBO}}(\phi,\psi)
= \E_{z \sim q_\phi}[\log p_\psi(\mathbf{x}|z)]
  - \KL{q_\phi(z|\mathbf{x})}{p(z)}.
\]
\end{definition}

\subsection{Game Module: Mean-Field Equilibrium Operator}

The core innovation of our architecture is replacing the recurrent world
model $M$ with a mean-field equilibrium operator.

\begin{definition}[Market mean-field game]
\label{def:game}
The market consists of a continuum of agents $u \in [0,1]$. Agent $u$'s
latent state evolves as:
\begin{equation}
\label{eq:agent_sde}
dz^u_t = f(z^u_t, a^u_t, \mu_t, u)\,dt + \sigma\,dW^u_t,
\quad z^u_0 \sim \rho_0,
\end{equation}
where $a^u_t \in \A$ is the action, $\mu_t = \Law(z^u_t)$ is the
\emph{mean field} (distribution of all agents), and $(W^u)_{u\in[0,1]}$
are independent Brownian motions.  Agent $u$ minimises the cost:
\begin{equation}
\label{eq:agent_cost}
J^u(a) = \E\!\left[\int_0^T \ell(z^u_t, a^u_t, \mu_t, u)\,dt
                   + g(z^u_T, \mu_T, u)\right].
\end{equation}
\end{definition}

\begin{theorem}[MFG equilibrium existence]
\label{thm:mfg_exist}
Under the following conditions:
\begin{enumerate}[label=(\textnormal{A\arabic*})]
  \item $f$ and $\ell$ are jointly continuous in all arguments and
        Lipschitz in $(z,a,\mu)$ with constant $L_f$;
  \item The action space $\A$ is compact and convex;
  \item There exists a unique minimiser $\hat{a}(z,\mu,u) = \argmin_{a\in\A}\mathcal{H}(z,a,\mu,u,p)$
        for the Hamiltonian $\mathcal{H}$ with adjoint $p$;
  \item (\emph{Lasry--Lions monotonicity})
        $\int[\ell(z,\mu,a)-\ell(z,\mu',a)]\,d(\mu-\mu')(z) \geq 0$
        for all $\mu, \mu', a$,
\end{enumerate}
there exists a unique mean-field game equilibrium
$(\hat{a}^*, \hat{\mu})$ satisfying:
(i)~$\hat{a}^*_t = \argmin_{a}J^u(a|\hat{\mu})$ for $\mu$-a.e.\ $u$;
(ii)~$\hat{\mu}_t = \Law(z^u_t|\hat{a}^*)$ for all $t$.
\end{theorem}

\begin{proof}
Define the fixed-point map $\Phi: C([0,T];\Prob_2(\Z)) \to C([0,T];\Prob_2(\Z))$
by: given a flow $\mathbf{m} = (\mu_t)$, solve the forward SDE~\eqref{eq:agent_sde}
with optimal feedback $\hat{a}(z,\mu_t,u)$ and output $\Phi(\mathbf{m})_t = \Law(z^u_t)$.

\textit{Self-mapping.} By (A1)--(A2) and Gronwall, solutions of~\eqref{eq:agent_sde}
satisfy a uniform $L^2$ bound, so $\Phi(\mathbf{m}) \in C([0,T];\Prob_2(\Z))$.

\textit{Compactness.} Tightness of $\{\Phi(\mathbf{m}_n)\}$ for bounded sequences
$\{\mathbf{m}_n\}$ follows from Arzela--Ascoli applied to the SDE solutions.

\textit{Continuity.} If $\mathbf{m}_n \to \mathbf{m}$ in $C([0,T];\Prob_2(\Z))$
(Wasserstein-2 topology), then by (A1) and stability of SDEs,
$\Phi(\mathbf{m}_n) \to \Phi(\mathbf{m})$.

\textit{Schauder's theorem.} $\Phi$ maps the closed convex set
$\mathcal{K} = \{\mathbf{m}: \sup_t\WD{\mu_t}{\rho_0}^2 \leq R\}$
(for $R$ large enough) into itself, continuously, and $\mathcal{K}$
is compact. Schauder's fixed-point theorem~\citep{Evans2010} yields
a fixed point $\hat{\mathbf{m}}$ with $\Phi(\hat{\mathbf{m}}) = \hat{\mathbf{m}}$,
i.e., a MFG equilibrium.

\textit{Uniqueness.} The Lasry--Lions condition (A4) ensures that any
two equilibria $(\hat{a}^*, \hat{\mu})$ and $(\hat{a}^{**},\hat{\mu}^{**})$ satisfy
$\WD{\hat{\mu}_t}{\hat{\mu}^{**}_t} = 0$ for all $t$ (by a standard
energy estimate using the monotonicity of $\ell$), hence uniqueness.
\end{proof}

\subsection{Controller Module C}

\begin{definition}[Optimal controller]
The controller $\pi^*: \Z \times \Hil_\mathrm{RNN} \to \A$ maps the current
latent state $z_t$ and world-model hidden state $h_t$ to an optimal action:
\[
\pi^*(z_t, h_t) = \argmax_{a\in\A}\, Q^*(z_t, h_t, a),
\]
where $Q^*$ is the optimal action-value function satisfying the HJB equation:
\[
0 = \partial_t Q + \sup_{a\in\A}\left[
    f(z,a,\hat{\mu}_t)^\top \nabla_z Q + \ell(z,a,\hat{\mu}_t)
  \right]
  + \tfrac{1}{2}\sigma^2\Delta_z Q.
\]
\end{definition}

% ============================================================
\section{Hierarchical Mean-Field Theory}
\label{sec:mfg}
% ============================================================

\subsection{Motivation: Multi-Scale Structure of Financial Markets}

A key structural observation motivates the hierarchical extension:

\begin{enumerate}
\item At the \textbf{national/regional scale}, sovereign entities (countries,
  central banks, regulatory bodies) interact without a supranational central
  planner—a MFG structure.
\item Within each nation, individual firms are subject to regulation and
  monetary policy set by the central bank—a partial MFC structure.
\item At the \textbf{firm level}, individual investors within a firm's investor
  base again interact without coordination—a MFG structure.
\end{enumerate}

This hierarchical interplay—MFG at the macro level, mixed MFG/MFC at the
meso level—cannot be captured by a single-level mean-field model.

\subsection{Two-Level McKean--Vlasov System}

Let $N \in \N$ denote the number of groups (nations) and $M$ the size of the
continuum of agents within each group.

\paragraph{Level 1 (macro-group dynamics).}
Each group $j \in \{1,\ldots,N\}$ has a representative state
$\xi^j_t \in \R^{d_1}$ evolving as:
\begin{equation}
\label{eq:level1_sde}
d\xi^j_t = b_1(\xi^j_t, \nu^{(1)}_t, \alpha^j_t)\,dt + \sigma_1\,dB^j_t,
\quad \xi^j_0 = \xi^j_0 \in \R^{d_1},
\end{equation}
where $\nu^{(1)}_t = \frac{1}{N}\sum_{k=1}^N \delta_{\xi^k_t}$ is the
empirical measure of macro-group states, $\alpha^j_t$ is the policy
(control) of group $j$, and $(B^j)_{j=1}^N$ are independent Brownian motions.
Group $j$ minimises the macro-cost:
\begin{equation}
\label{eq:level1_cost}
J_1^j(\alpha^j) = \E\!\left[\int_0^T f_1(\xi^j_t, \nu^{(1)}_t, \alpha^j_t)\,dt
                             + g_1(\xi^j_T, \nu^{(1)}_T)\right].
\end{equation}

\paragraph{Level 2 (micro-agent dynamics).}
Within group $j$, a continuum of agents $u \in [0,1]$ has states
$x^{j,u}_t \in \R^{d_2}$ evolving as:
\begin{equation}
\label{eq:level2_sde}
dx^{j,u}_t = b_2(x^{j,u}_t, \mu^{(2)}_{j,t}, \xi^j_t, a^{j,u}_t)\,dt
             + \sigma_2\,dW^{j,u}_t,
\end{equation}
where $\mu^{(2)}_{j,t} = \int_0^1 \delta_{x^{j,u}_t}\,du$ is the intra-group
mean field, $a^{j,u}_t$ is agent $u$'s control, and $(W^{j,u})$ are i.i.d.
Each agent minimises:
\begin{equation}
\label{eq:level2_cost}
J_2^{j,u}(a) = \E\!\left[\int_0^T f_2(x^{j,u}_t, \mu^{(2)}_{j,t}, \xi^j_t, a^{j,u}_t)\,dt
                          + g_2(x^{j,u}_T, \mu^{(2)}_{j,T}, \xi^j_T)\right].
\end{equation}

\paragraph{Coupling condition.}
The Level-1 and Level-2 dynamics are coupled through the \emph{aggregation
functional}
\begin{equation}
\label{eq:agg}
\Psi_j\!\left(\mu^{(2)}_{j,\cdot}\right)
:= \int_{\R^{d_2}} \phi(x)\,\mu^{(2)}_{j,t}(dx), \quad t \in [0,T],
\end{equation}
for a bounded Lipschitz test function $\phi$; the macro-group cost $J_1^j$
depends on $\Psi_j(\mu^{(2)}_{j,\cdot})$ through the drift $b_1$.

\subsection{Hierarchical Nash Equilibrium}

\begin{definition}[Hierarchical Nash equilibrium]
\label{def:hier_nash}
A pair $\bigl((\alpha^{j,*})_{j=1}^N,\, (a^{j,u,*})_{j,u}\bigr)$ is a
\emph{hierarchical Nash equilibrium} if:
\begin{enumerate}[label=(\roman*)]
\item For each group $j$ and the resulting Level-2 equilibrium
  $\hat{\mu}^{(2)}_j[\xi^j]$ (defined in Definition~\ref{def:lvl2_eq} below),
  the macro policy satisfies:
  \[
  J_1^j(\alpha^{j,*}) \leq J_1^j(\alpha^j)
  \quad \forall\,\alpha^j,\, \text{given } (\alpha^{k,*})_{k\neq j};
  \]
\item For each group $j$ and each agent $u \in [0,1]$, given $\xi^j$ and
  $\mu^{(2)}_{j,t} = \hat{\mu}^{(2)}_{j,t}[\xi^j]$:
  \[
  J_2^{j,u}(a^{j,u,*}) \leq J_2^{j,u}(a^{j,u}) \quad \forall\, a^{j,u}.
  \]
\end{enumerate}
\end{definition}

\subsection{Assumptions}

\begin{assumption}[Regularity at Level 2]
\label{ass:L2}
\begin{enumerate}[label=(L\arabic*)]
\item $b_2$, $f_2$, $g_2$ are jointly continuous and Lipschitz in
  $(x, \mu, \xi, a)$ with constant $L_2 < \infty$, uniformly in $j$.
\item $b_2$ satisfies linear growth: $|b_2(x,\mu,\xi,a)| \leq C_2(1+|x|)$.
\item The action space $\A_2$ is compact and convex.
\item Lasry--Lions monotonicity at Level 2:
  $\int[f_2(x,\mu,\xi,a) - f_2(x,\mu',\xi,a)]\,d(\mu-\mu')(x) \geq 0$.
\end{enumerate}
\end{assumption}

\begin{assumption}[Regularity at Level 1]
\label{ass:L1}
\begin{enumerate}[label=(L\arabic*), start=5]
\item $b_1$, $f_1$, $g_1$ are jointly continuous and Lipschitz in
  $(\xi, \nu, \alpha)$ with constant $L_1 < \infty$.
\item The action space $\A_1$ is compact and convex.
\item Lasry--Lions monotonicity at Level 1:
  $\int[f_1(\xi,\nu,\alpha) - f_1(\xi,\nu',\alpha)]\,d(\nu-\nu')(\xi) \geq 0$.
\item The aggregation map $\Psi_j: C([0,T];\Prob_2(\R^{d_2})) \to C([0,T];\R)$
  is Lipschitz with constant $L_\Psi < \infty$.
\end{enumerate}
\end{assumption}

\subsection{Main Theorem}

\begin{theorem}[Hierarchical Mean-Field Nash Equilibrium]
\label{thm:hier_mfg}
Under Assumptions~\ref{ass:L2} and~\ref{ass:L1}, there exists a unique
hierarchical Nash equilibrium in the sense of Definition~\ref{def:hier_nash}.
\end{theorem}

The proof proceeds through a sequence of four lemmas.

\begin{definition}[Level-2 equilibrium operator]
\label{def:lvl2_eq}
For each fixed continuous path $\xi = (\xi_t)_{t\in[0,T]} \in
C([0,T];\R^{d_1})$, let $\mathrm{MFG}_2[\xi]$ denote the mean-field game
with dynamics~\eqref{eq:level2_sde}--\eqref{eq:level2_cost} with $\xi^j_t$
replaced by $\xi_t$.  By Theorem~\ref{thm:mfg_exist} (applied to the
Level-2 system under (L1)--(L4)), $\mathrm{MFG}_2[\xi]$ has a unique
equilibrium measure flow $\hat{\mu}^{(2)}[\xi] \in C([0,T];\Prob_2(\R^{d_2}))$.
\end{definition}

\begin{lemma}[Level-2 stability estimate]
\label{lem:L2_stability}
Under Assumption~\ref{ass:L2}, for any two paths
$\xi, \xi' \in C([0,T];\R^{d_1})$:
\[
\sup_{t\in[0,T]} \WD{\hat{\mu}^{(2)}_t[\xi]}{\hat{\mu}^{(2)}_t[\xi']}
\leq C_{2,\mathrm{stab}}\,\norm{\xi - \xi'}_{C([0,T])},
\]
where $C_{2,\mathrm{stab}} = C_{2,\mathrm{stab}}(L_2, T, \sigma_2) < \infty$
depends only on the problem parameters.
\end{lemma}

\begin{proof}
Let $x_t$ and $x'_t$ be the optimal state processes for $\mathrm{MFG}_2[\xi]$
and $\mathrm{MFG}_2[\xi']$ respectively.  Define $\Delta_t = x_t - x'_t$.
By the Lipschitz condition (L1), the dynamics of $\Delta_t$ satisfy:
\[
d\Delta_t = \left[b_2(x_t, \hat\mu_t, \xi_t, \hat a_t)
                 - b_2(x'_t, \hat\mu'_t, \xi'_t, \hat a'_t)\right]dt.
\]
Using the triangle inequality for the Lipschitz bound:
\begin{align*}
|d\Delta_t| &\leq L_2\bigl(|\Delta_t| + \WD{\hat\mu_t}{\hat\mu'_t}
                           + |\xi_t - \xi'_t| + |\hat a_t - \hat a'_t|\bigr)dt.
\end{align*}
Since the optimal control satisfies the same Lipschitz bound (by the
smoothness of the Hamiltonian), and $\WD{\hat\mu_t}{\hat\mu'_t}
\leq (\E[|\Delta_t|^2])^{1/2}$, Gronwall's inequality gives:
\[
\E\!\left[\sup_{t\leq T}|\Delta_t|^2\right]
\leq C e^{C L_2 T}\norm{\xi - \xi'}_{C([0,T])}^2.
\]
Since $\hat\mu^{(2)}_t[\xi] = \Law(x_t)$ and $\hat\mu^{(2)}_t[\xi'] = \Law(x'_t)$,
the $W_2$ distance is bounded by $(\E[|\Delta_t|^2])^{1/2}$, yielding the claim.
\end{proof}

\begin{lemma}[Continuity of the aggregation functional]
\label{lem:agg_cont}
The composed map
$\xi \mapsto \Psi_j(\hat\mu^{(2)}[\xi])$
from $C([0,T];\R^{d_1})$ to $C([0,T];\R)$ is Lipschitz with constant
$L_\Psi \cdot C_{2,\mathrm{stab}}$.
\end{lemma}

\begin{proof}
Direct application of (L8) and Lemma~\ref{lem:L2_stability}:
\begin{align*}
\norm{\Psi_j(\hat\mu^{(2)}[\xi]) - \Psi_j(\hat\mu^{(2)}[\xi'])}_{C([0,T])}
&\leq L_\Psi \sup_t \WD{\hat\mu^{(2)}_t[\xi]}{\hat\mu^{(2)}_t[\xi']} \\
&\leq L_\Psi C_{2,\mathrm{stab}}\norm{\xi - \xi'}_{C([0,T])}.
\end{align*}
\end{proof}

\begin{lemma}[Existence of Level-1 MFG equilibrium]
\label{lem:L1_exist}
Fix the Level-2 equilibrium measure flows $\hat\mu^{(2)}[\xi^j]$ as a function
of the Level-1 paths $\xi^j$. Then the Level-1 system~\eqref{eq:level1_sde}--\eqref{eq:level1_cost} (with the aggregation term incorporated into $b_1$) admits a unique MFG equilibrium $(\hat\alpha^{j,*}, \hat\nu^{(1)})$.
\end{lemma}

\begin{proof}
The Level-1 system is a standard mean-field game with state space $\R^{d_1}$
and action space $\A_1$.  By Lemma~\ref{lem:agg_cont}, incorporating
$\Psi_j(\hat\mu^{(2)}[\xi^j])$ into the Level-1 drift $b_1$ produces a new
drift $\tilde{b}_1(\xi^j, \nu^{(1)}, \alpha^j) := b_1(\xi^j, \nu^{(1)}, \alpha^j,
\Psi_j(\hat\mu^{(2)}[\xi^j]))$ that is Lipschitz in $(\xi^j, \nu^{(1)}, \alpha^j)$
with constant $\max(L_1, L_\Psi C_{2,\mathrm{stab}})$.
Assumptions (L5)--(L7) ensure that Theorem~\ref{thm:mfg_exist} applies to the
Level-1 system with this modified drift, giving the existence and uniqueness
of $(\hat\alpha^{j,*}, \hat\nu^{(1)})$.
\end{proof}

\begin{lemma}[Consistency: Level-2 under optimal Level-1 policies]
\label{lem:consistency}
Given the Level-1 equilibrium trajectories $(\hat\xi^j_t)$ generated by
$\hat\alpha^{j,*}$, the corresponding Level-2 equilibria
$\hat{a}^{j,u,*} = \hat{a}^{j,u,*}[\hat\xi^j]$ satisfy the Level-2
optimality conditions in Definition~\ref{def:hier_nash}(ii).
\end{lemma}

\begin{proof}
By Definition~\ref{def:lvl2_eq}, $\hat{a}^{j,u,*}[\hat\xi^j]$ is the
equilibrium of $\mathrm{MFG}_2[\hat\xi^j]$, which directly gives optimality.
No further argument is needed beyond the definition.
\end{proof}

\begin{proof}[Proof of Theorem~\ref{thm:hier_mfg}]
\textit{Existence.}
\begin{enumerate}
\item By Theorem~\ref{thm:mfg_exist} applied at Level 2 (Lemmas exist for any
  fixed $\xi^j$), the Level-2 equilibrium operator $\xi^j \mapsto
  (\hat{a}^{j,u,*}[\xi^j], \hat\mu^{(2)}_j[\xi^j])$ is well-defined.
\item By Lemma~\ref{lem:L2_stability} and~\ref{lem:agg_cont}, the aggregation
  functional $\xi^j \mapsto \Psi_j(\hat\mu^{(2)}_j[\xi^j])$ is Lipschitz.
\item By Lemma~\ref{lem:L1_exist}, the Level-1 MFG with the aggregated
  feedback has a unique equilibrium $(\hat\alpha^{j,*}, \hat\nu^{(1)})$,
  generating equilibrium trajectories $\hat\xi^j_t$.
\item By Lemma~\ref{lem:consistency}, the Level-2 policies
  $\hat{a}^{j,u,*}[\hat\xi^j]$ are optimal given $\hat\xi^j$.
  Together with the Level-1 optimality from step 3, Definition~\ref{def:hier_nash}
  is satisfied.
\end{enumerate}

\textit{Uniqueness.}
Suppose $\bigl((\alpha^{j,*}),\, (a^{j,u,*})\bigr)$ and
$\bigl((\alpha^{j,**}),\, (a^{j,u,**})\bigr)$ are two hierarchical Nash
equilibria generating trajectory distributions $\nu^{(1)}$ and $\nu^{(1),**}$
respectively.

At Level 2: by the uniqueness part of Theorem~\ref{thm:mfg_exist} (uniqueness
follows from the Lasry--Lions condition (L4)), for each fixed Level-1
trajectory path $\xi^j$:
$\hat\mu^{(2)}_j[\xi^j]$ is unique, so the Level-2 equilibria coincide whenever
the Level-1 paths coincide.

At Level 1: the Level-1 system with the (unique) aggregation
$\Psi_j(\hat\mu^{(2)}_j[\xi^j])$ is a standard MFG satisfying (L5)--(L7),
whose unique equilibrium measure flow $\hat\nu^{(1)}$ is guaranteed by the
Lasry--Lions condition (L7).  Hence $\nu^{(1)} = \nu^{(1),**}$, and the
Level-1 equilibrium trajectories coincide.  Feeding these back into the
Level-2 uniqueness gives $a^{j,u,*} = a^{j,u,**}$ $\mu_j$-a.e.
\end{proof}

\begin{corollary}[Multi-scale information decomposition]
\label{cor:multiscale}
Under the hierarchical equilibrium, the optimal micro-agent policy decomposes:
\[
a^{j,u,*}_t
= \underbrace{\vp_1(\hat\xi^j_t, \hat\nu^{(1)}_t)}_{\text{macro contribution}}
+ \underbrace{\vp_2(x^{j,u}_t, \hat\mu^{(2)}_{j,t}, \hat\xi^j_t)}_{\text{micro contribution}}
+ O(\epsilon),
\]
where $\epsilon = L_\Psi C_{2,\mathrm{stab}} \cdot \sup_t|\xi^j_t - \hat\xi^j_t|$
measures the coupling strength.
\end{corollary}

\begin{proof}
Expand the optimal value function $V^{j,u}(x, \mu, \xi)$ around the equilibrium
$(\hat\mu^{(2)}_{j,t}, \hat\xi^j_t)$ using Taylor's theorem. The first-order
terms in $\xi$ yield $\vp_1$ via the envelope theorem, and the terms in $\mu$
yield $\vp_2$. The $O(\epsilon)$ remainder follows from Lemma~\ref{lem:L2_stability}.
\end{proof}

% ============================================================
\section{Stochastic Lyapunov Stability}
\label{sec:stability}
% ============================================================

\subsection{Lyapunov Stability under Dual Noise}

The classical Lyapunov stability theory for ODEs must be substantially extended
when the noise includes a jump component. The key tool is the Itô formula for
Lévy-driven processes.

\begin{definition}[Infinitesimal generator]
\label{def:generator}
For the process $\mathbf{S}_t$ satisfying~\eqref{eq:full_sde} (ignoring
event jumps for clarity), the \emph{extended generator} acting on
$C^2(\Hil)$ is:
\begin{align}
\label{eq:generator}
\mathcal{L}V(\mathbf{S})
&= \nabla V(\mathbf{S})^\top \mu(\mathbf{S})
   + \tfrac{1}{2}\tr\!\left(\sigma_\tau\sigma_\tau^\top \nabla^2 V(\mathbf{S})\right) \nonumber\\
&\quad + \int_{\R^k}\!\left[V(\mathbf{S}+z) - V(\mathbf{S})
         - z^\top\nabla V(\mathbf{S})\,\1_{|z|\leq 1}\right]\nu^\eta(dz).
\end{align}
\end{definition}

The three terms correspond respectively to: the contribution of the
deterministic drift, the diffusive physical noise, and the behavioral jump noise.

\begin{theorem}[Stochastic Lyapunov stability and invariant measure]
\label{thm:lyapunov}
Suppose there exists a $C^2$ function $V: \Hil \to \R_+$ and constants
$c, c_0 > 0$ such that:
\begin{equation}
\label{eq:lyapunov_cond}
\mathcal{L}V(\mathbf{S}) \leq -c\,V(\mathbf{S}) + c_0
\quad \text{for all } \mathbf{S} \in \Hil.
\end{equation}
Then:
\begin{enumerate}[label=(\roman*)]
\item \textbf{Exponential ergodicity}: For any initial condition $\mathbf{S}_0$,
\[
\E[V(\mathbf{S}_t)] \leq V(\mathbf{S}_0)\,e^{-ct} + \frac{c_0}{c}.
\]
\item \textbf{Existence of invariant measure}: The process $\mathbf{S}_t$ admits
at least one invariant (stationary) probability measure $\pi^*$.
\item \textbf{Exponential convergence}: If additionally $V \geq 1$ and the
sublevel sets $\{V \leq r\}$ are compact (petite sets for the Markov
chain skeleton), then $\pi^*$ is unique and
\[
\norm{\PP(\mathbf{S}_t \in \cdot|\mathbf{S}_0) - \pi^*}_{\mathrm{TV}}
\leq C(\mathbf{S}_0)\,e^{-\rho t}
\]
for some $C(\mathbf{S}_0) < \infty$ and $\rho > 0$ depending on $c$ and $c_0$.
\end{enumerate}
\end{theorem}

\begin{proof}
\textit{(i) Exponential bound.}
By Itô's formula for Lévy processes~\citep{Sato1999,Protter2005}:
\begin{align*}
dV(\mathbf{S}_t) &= \mathcal{L}V(\mathbf{S}_t)\,dt
 + \nabla V(\mathbf{S}_t)^\top\sigma_\tau\,dW_t \\
&\quad + \int_{\R^k}\!\bigl[V(\mathbf{S}_{t^-}+z) - V(\mathbf{S}_{t^-})\bigr]\,
         \tilde{N}^\eta(dt,dz).
\end{align*}
The stochastic integrals are local martingales. Taking expectations:
\[
\frac{d}{dt}\E[V(\mathbf{S}_t)]
= \E[\mathcal{L}V(\mathbf{S}_t)]
\leq -c\,\E[V(\mathbf{S}_t)] + c_0.
\]
Gronwall's inequality gives $\E[V(\mathbf{S}_t)] \leq V(\mathbf{S}_0)e^{-ct} + c_0/c$.

\textit{(ii) Existence of invariant measure.}
The bound in (i) shows that the family $\{(1/t)\int_0^t \PP(\mathbf{S}_s \in \cdot)\,ds\}_{t>0}$
is tight in $\Prob(\Hil)$ (by Markov's inequality and the upper bound
$c_0/c$ on $\E[V]$). Any weak limit point of this family is an invariant measure
(Krylov--Bogoliubov theorem~\citep{DaPratoZabczyk1996}).

\textit{(iii) Uniqueness and exponential rate.}
Under the petite-set condition, the process is \emph{$V$-geometrically ergodic}
in the sense of~\citet{MeynTweedie1993} (their Foster--Lyapunov
Criterion V). The uniqueness of $\pi^*$ and the TV convergence follow from
their Theorem 15.0.1.
\end{proof}

\begin{remark}[Lyapunov condition in financial terms]
The condition~\eqref{eq:lyapunov_cond} with $V(\mathbf{S}) = \|\mathbf{S} - \mathbf{S}^*\|^2$
reduces to a mean-reversion condition: the drift and jump terms together must
push the market state back toward $\mathbf{S}^*$ at rate at least $c$, i.e.,
\[
2(\mathbf{S}-\mathbf{S}^*)^\top\mu(\mathbf{S})
+ \tr(\sigma_\tau\sigma_\tau^\top)
+ \int_{\R^k}\left(2(\mathbf{S}-\mathbf{S}^*)^\top z + |z|^2\right)\nu^\eta(dz)
\leq -c\|\mathbf{S}-\mathbf{S}^*\|^2 + c_0.
\]
\end{remark}

\subsection{Bubbles and Crises as Lyapunov Violations}

\begin{definition}[Instability region]
The \emph{instability region} is $\mathcal{R}_{\mathrm{unstable}}
:= \{\mathbf{S} \in \Hil : \mathcal{L}V(\mathbf{S}) > 0\}$.
\end{definition}

\begin{proposition}[Analytic signature of market bubbles]
\label{prop:bubble}
If $\mathbf{S}_0 \in \mathcal{R}_{\mathrm{unstable}}$ and the behavioral jump
intensity $\lambda_\eta(t)$ is sufficiently large (high-sentiment regime), then:
\[
\PP\!\left(\tau_{\mathrm{escape}} < T\right) > 0,
\quad \tau_{\mathrm{escape}} := \inf\{t > 0 : V(\mathbf{S}_t) > R\}
\]
for any $R > V(\mathbf{S}_0)$.  That is, the market state escapes to a
high-$V$ region (bubble or crash trajectory) with positive probability.
\end{proposition}

\begin{proof}
In $\mathcal{R}_{\mathrm{unstable}}$, $\mathcal{L}V > 0$, so
$\E[V(\mathbf{S}_{t\wedge\tau})]$ is locally increasing. By the optional
stopping theorem applied to the stopped submartingale
$V(\mathbf{S}_{t\wedge\tau})$, the probability of exiting $\{V \leq R\}$
before time $T$ is bounded below by
$\PP(\tau_{\mathrm{escape}} < T) \geq
\mathcal{L}V(\mathbf{S}_0)/(R - V(\mathbf{S}_0)) \cdot \delta + O(\delta^2)$
for small $\delta > 0$ (first-exit probability estimate).
\end{proof}

\begin{corollary}[Real-time market risk indicator]
Define the instantaneous risk indicator:
\[
\mathrm{RI}(t) := \mathcal{L}V(\mathbf{S}_t).
\]
$\mathrm{RI}(t) > 0$ signals that the market is in a potential bubble
or crisis regime; $\mathrm{RI}(t) < -c\,V(\mathbf{S}_t)$ signals strong
mean-reversion. Monitoring $\mathrm{RI}$ provides a theoretically
grounded, model-consistent early-warning indicator.
\end{corollary}

% ============================================================
\section{Unified Evolution Equation}
\label{sec:unified}
% ============================================================

We now assemble all components into the complete evolution equation.

\begin{theorem}[Unified market world model equation]
\label{thm:unified}
The market state process $\mathbf{S}_t$ under the E-Game-C architecture,
with dual noise and financial event operators, satisfies:
\begin{equation}
\label{eq:unified}
\mathbf{S}_t = \mathbf{S}_0
  + \int_0^t \mu^*(\mathbf{S}_u, \hat\mu_u)\,du
  + \int_0^t \sigma_\tau(\mathbf{S}_u)\,dW^{(\tau)}_u
  + \int_0^t\!\!\int_{\R^k} z\,\tilde{N}^\eta(du,dz)
  + \sum_{\substack{w \in \mathcal{W}_I \cup \mathcal{W}_{II} \\ \tau_w \leq t}}
    (T_w - I)\mathbf{S}_{\tau_w^-}
  + \sum_{\substack{w \in \mathcal{W}_{III} \\ \tau_w \leq t}}
    \mathcal{R}_w(\mathbf{S}_{\tau_w^-}),
\end{equation}
where $\mu^*(\mathbf{S}, \hat\mu)$ is the equilibrium drift from the
hierarchical MFG (Theorem~\ref{thm:hier_mfg}), $T_w$ are Type-I/II event
operators, and $\mathcal{R}_w$ are Type-III (dimension-changing) groupoid
morphisms.
\end{theorem}

\begin{proof}
The equation follows by assembling:
(a) the Itô integral form (Definition~\ref{def:ito_integral});
(b) the dual noise decomposition (Theorem~\ref{thm:noise_decomp});
(c) the event operator jumps (Definition~\ref{def:operator});
(d) the equilibrium drift from the hierarchical MFG equilibrium
(Theorem~\ref{thm:hier_mfg}), which replaces the generic $\mu(\mathbf{S},u)$.
Well-posedness of~\eqref{eq:unified} follows from the Lipschitz conditions in
Assumptions~\ref{ass:L2}--\ref{ass:L1} and the standard existence theory for
McKean--Vlasov SDEs with jumps~\citep{MisierMeleard1996}.
\end{proof}

% ============================================================
\section{Discussion}
\label{sec:discussion}
% ============================================================

\subsection{Relation to Prior Work}

\paragraph{Factor models.}
Our framework subsumes factor models as a special case: setting $\mu^* = \sum_k \beta_k F_k$ (linear factor structure), $\sigma_\tau = \sigma_0 I$, $\nu^\eta = 0$, and $T_w = I$ (no events) recovers a multi-factor APT model. The advantage of our setting is that it does not assume stationarity of factor loadings.

\paragraph{MFG in finance.}
Existing MFG-based financial models~\citep{CarmonaLacker2015,CarmonaMouzouni2023}
operate at a single scale.  Our hierarchical extension (Theorem~\ref{thm:hier_mfg})
is, to our knowledge, the first to prove existence of a hierarchical Nash
equilibrium in a two-level McKean--Vlasov system with explicit Level-1/Level-2
coupling via an aggregation functional.

\paragraph{World models.}
The world-model paradigm has been applied to portfolio
optimisation~\citep{AwmlgLiverpool2026} and order-book
simulation~\citep{MorganJPMaGAN} but without a theoretical foundation.
Our E-Game-C architecture provides this foundation and subsumes prior
architectures as special cases (e.g., \citet{HaSchmidhuber2018} corresponds to
a Level-1-only model with no behavioral noise).

\subsection{Limitations and Future Work}

\paragraph{Empirical validation.}
This paper establishes the theoretical framework. Empirical validation—
calibrating $(\sigma_\tau, \nu^\eta)$ from high-frequency data, estimating
$\hat\mu^{(2)}$ from a public dataset, testing the Lyapunov risk indicator—
is deferred to subsequent work.

\paragraph{Graphon extension.}
The agent type $u$ is currently a scalar in $[0,1]$.  As noted in
Section~\ref{sec:mfg}, the natural extension is $u \in [0,1]^m$ (a vector of
agent characteristics), leading to Graphon mean-field games~\citep{CarmonaGraphon2022}.
This extension is theoretically straightforward given our framework but requires
additional functional-analytic machinery.

\paragraph{Non-Markovian extensions.}
The Markovian assumption (McKean--Vlasov SDE) can be relaxed to allow
path-dependent coefficients, at the cost of working in infinite-dimensional
path spaces.

% ============================================================
\section{Conclusion}
\label{sec:conclusion}
% ============================================================

We have presented the first systematic mathematical theory for market world
models. Starting from the observation that financial markets have no correct
ODE description (Proposition~\ref{prop:nondiff}), we developed:

\begin{itemize}
\item A dual noise decomposition into physical (Brownian) and behavioral (Lévy-jump)
  components, with a Cramér--Rao lower bound (Theorem~\ref{thm:cramer_rao})
  that quantifies the irreducible prediction uncertainty of any world model.

\item A groupoid algebra of financial event operators
  (Theorem~\ref{thm:groupoid}), providing the first algebraically rigorous
  treatment of corporate events as structured transformations of market state.

\item A hierarchical mean-field theory with a complete existence and uniqueness
  proof of the hierarchical Nash equilibrium (Theorem~\ref{thm:hier_mfg}),
  capturing the simultaneous presence of MFG and MFC structures at different
  market scales.

\item A stochastic Lyapunov stability theorem under combined Brownian-Lévy noise
  (Theorem~\ref{thm:lyapunov}), yielding a theoretically grounded real-time
  market risk indicator.
\end{itemize}

The E-Game-C architecture ties these components into an executable framework
whose equilibrium drift is provided by the hierarchical MFG and whose noise
structure is governed by the dual decomposition.

Together, these results place the world-model paradigm for quantitative
finance on a rigorous mathematical foundation comparable to the classical
Itô-calculus foundations of derivative pricing.

% ============================================================
\appendix
% ============================================================

\section{Technical Lemmas}
\label{app:lemmas}

\begin{lemma}[Gronwall's inequality, integral form]
\label{lem:gronwall}
If $u(t) \leq \alpha(t) + \int_0^t \beta(s)\,u(s)\,ds$ with $\alpha$ non-decreasing
and $\beta \geq 0$, then $u(t) \leq \alpha(t)\exp\!\int_0^t\beta(s)\,ds$.
\end{lemma}

\begin{lemma}[$W_2$ distance and optimal coupling]
\label{lem:w2}
For probability measures $\mu, \mu'$ on $\R^d$,
$W_2(\mu,\mu')^2 = \inf_{\gamma \in \Gamma(\mu,\mu')} \int |x-y|^2\,d\gamma(x,y)$,
and for measures arising as laws of random variables $X, X'$:
$W_2(\Law(X), \Law(X')) \leq (\E[|X-X'|^2])^{1/2}$.
\end{lemma}

\section{Proof of Non-Commutativity of Event Sequences}
\label{app:noncomm}

\begin{proposition}[Event sequences are non-commutative]
In general, $T_{w_2} \circ T_{w_1} \neq T_{w_1} \circ T_{w_2}$.
\end{proposition}

\begin{proof}
Take the single-asset state $s = (p, \kappa)^\top \in \R^2$ (log-price and
log-shares). Let $w_1$ be a 2-for-1 stock split: $T_{w_1}(s) = (p, \kappa + \log 2)^\top$.
Let $w_2$ be a cash dividend of $d$ per share:
$T_{w_2}(s) = (p - \log(1 + d/e^p), \kappa)^\top$ (ex-dividend price drop).

Then:
$(T_{w_2} \circ T_{w_1})(s) = T_{w_2}(p, \kappa+\log 2) = (p - \log(1+d/e^p), \kappa+\log 2)$,
while
$(T_{w_1} \circ T_{w_2})(s) = T_{w_1}(p-\log(1+d/e^p), \kappa) = (p-\log(1+d/e^p), \kappa+\log 2)$.
In this case they coincide. However, if $w_2$ is instead a rights issue that
changes $\kappa$: $T_{w_2}(s) = (p\cdot r^{-1}, \kappa + \log r)^\top$, then
$T_{w_2} \circ T_{w_1}(s) = T_{w_2}(p, \kappa+\log 2) = (p/r, \kappa+\log 2 + \log r)$,
while $T_{w_1} \circ T_{w_2}(s) = T_{w_1}(p/r, \kappa+\log r) = (p/r, \kappa+\log r + \log 2)$.
Since $\log 2 + \log r = \log r + \log 2$ the $\kappa$-coordinates agree, but the
price adjustment for the rights issue depends on the \emph{pre-split} price in
$T_{w_2} \circ T_{w_1}$ and the \emph{post-split} price in $T_{w_1} \circ T_{w_2}$
when the rights issue price depends on the current market price—giving different
results.
\end{proof}

% ============================================================
% Bibliography
% ============================================================

\begin{thebibliography}{99}

\bibitem[Ait-Sahalia and Todorov(2009)]{AitSahaliaTodorov2009}
Ait-Sahalia, Y.\ and Todorov, V.\ (2009).
Estimating volatility in the presence of jumps.
\textit{Journal of Financial Economics} \textbf{93}(3):421--441.

\bibitem[Bensoussan et al.(2013)]{BensoussanYam2013}
Bensoussan, A., Chau, M.H.M., and Yam, P.\ (2013).
Mean-field approach for large-scale interbank risk management systems.
\textit{Finance and Stochastics} \textbf{19}(4):799--843.

\bibitem[Carmona et al.(2013)]{CarmonaDelarueLachapelle2013}
Carmona, R., Delarue, F., and Lachapelle, A.\ (2013).
Control of McKean--Vlasov dynamics versus mean field games.
\textit{Mathematics and Financial Economics} \textbf{7}(2):131--166.

\bibitem[Carmona and Delarue(2018)]{CarmonaDelarue2018book}
Carmona, R.\ and Delarue, F.\ (2018).
\textit{Probabilistic Theory of Mean Field Games with Applications I--II}.
Springer.

\bibitem[Carmona and Lacker(2015)]{CarmonaLacker2015}
Carmona, R.\ and Lacker, D.\ (2015).
A probabilistic weak formulation of mean field games and applications.
\textit{Annals of Applied Probability} \textbf{25}(3):1189--1231.

\bibitem[Carmona and Mouzouni(2023)]{CarmonaMouzouni2023}
Carmona, R.\ and Mouzouni, C.\ (2023).
Systemic risk and mean field games.
\textit{Finance and Stochastics} \textbf{27}:1--46.

\bibitem[Carmona et al.(2022)]{CarmonaGraphon2022}
Carmona, R., Cooney, D., Graves, C., and Laurière, M.\ (2022).
Stochastic graphon games: I. The static case.
\textit{Mathematics of Operations Research} \textbf{47}(1):750--778.

\bibitem[Carmona and Sun(2018)]{CarmonaSun2018}
Carmona, R.\ and Sun, P.\ (2018).
Mean field games with common noise.
\textit{Annals of Applied Probability} \textbf{26}(6):3560--3589.

\bibitem[Connes(1994)]{Connes1994}
Connes, A.\ (1994).
\textit{Noncommutative Geometry}.
Academic Press.

\bibitem[Cont and Bouchaud(2000)]{ContBouchaud2000}
Cont, R.\ and Bouchaud, J.-P.\ (2000).
Herd behavior and aggregate fluctuations in financial markets.
\textit{Macroeconomic Dynamics} \textbf{4}(2):170--196.

\bibitem[Cont and Tankov(2004)]{ContTankov2004}
Cont, R.\ and Tankov, P.\ (2004).
\textit{Financial Modelling with Jump Processes}.
Chapman \& Hall/CRC.

\bibitem[Da Prato and Zabczyk(1996)]{DaPratoZabczyk1996}
Da Prato, G.\ and Zabczyk, J.\ (1996).
\textit{Ergodicity for Infinite Dimensional Systems}.
Cambridge University Press.

\bibitem[Evans(2010)]{Evans2010}
Evans, L.C.\ (2010).
\textit{Partial Differential Equations}, 2nd ed.
American Mathematical Society.

\bibitem[Fama(1970)]{Fama1970}
Fama, E.F.\ (1970).
Efficient capital markets: A review of theory and empirical work.
\textit{Journal of Finance} \textbf{25}(2):383--417.

\bibitem[Ha and Schmidhuber(2018)]{HaSchmidhuber2018}
Ha, D.\ and Schmidhuber, J.\ (2018).
World models.
\textit{arXiv preprint arXiv:1803.10122}.

\bibitem[Hafner et al.(2020)]{Hafner2020Dreamer}
Hafner, D., Lillicrap, T., Ba, J., and Norouzi, M.\ (2020).
Dream to control: Learning behaviors by latent imagination.
\textit{International Conference on Learning Representations}.

\bibitem[Hafner et al.(2023)]{Hafner2023DreamerV3}
Hafner, D., Lillicrap, T., Norouzi, M., and Ba, J.\ (2023).
Mastering diverse domains through world models.
\textit{arXiv preprint arXiv:2301.04104}.

\bibitem[Huang et al.(2006)]{HuangMalhame2006}
Huang, M., Malhamé, R.P., and Caines, P.E.\ (2006).
Large population stochastic dynamic games: closed-loop McKean--Vlasov systems
and the Nash certainty equivalence principle.
\textit{Communications in Information and Systems} \textbf{6}(3):221--252.

\bibitem[Khasminskii(2012)]{Khasminskii2012}
Khasminskii, R.\ (2012).
\textit{Stochastic Stability of Differential Equations}, 2nd ed.
Springer.

\bibitem[Kingma and Welling(2014)]{KingmaWelling2014}
Kingma, D.P.\ and Welling, M.\ (2014).
Auto-encoding variational Bayes.
\textit{International Conference on Learning Representations}.

\bibitem[Kou(2002)]{Kou2002}
Kou, S.G.\ (2002).
A jump-diffusion model for option pricing.
\textit{Management Science} \textbf{48}(8):1086--1101.

\bibitem[Lasry and Lions(2007)]{LasryLions2007}
Lasry, J.-M.\ and Lions, P.-L.\ (2007).
Mean field games.
\textit{Japanese Journal of Mathematics} \textbf{2}(1):229--260.

\bibitem[Lux and Marchesi(1999)]{LuxMarchesi1999}
Lux, T.\ and Marchesi, M.\ (1999).
Scaling and criticality in a stochastic multi-agent model of a financial market.
\textit{Nature} \textbf{397}:498--500.

\bibitem[Meyn and Tweedie(1993)]{MeynTweedie1993}
Meyn, S.P.\ and Tweedie, R.L.\ (1993).
Stability of Markovian processes III: Foster--Lyapunov criteria for
continuous-time processes.
\textit{Advances in Applied Probability} \textbf{25}(3):518--548.

\bibitem[Méléard(1996)]{MisierMeleard1996}
Méléard, S.\ (1996).
Asymptotic behaviour of some interacting particle systems; McKean--Vlasov and
Boltzmann models.
\textit{Probabilistic Models for Nonlinear Partial Differential Equations},
Lecture Notes in Mathematics, Springer, 42--95.

\bibitem[Protter(2005)]{Protter2005}
Protter, P.E.\ (2005).
\textit{Stochastic Integration and Differential Equations}, 2nd ed.
Springer.

\bibitem[Ross(1976)]{Ross1976}
Ross, S.A.\ (1976).
The arbitrage theory of capital asset pricing.
\textit{Journal of Economic Theory} \textbf{13}(3):341--360.

\bibitem[Sato(1999)]{Sato1999}
Sato, K.-I.\ (1999).
\textit{Lévy Processes and Infinitely Divisible Distributions}.
Cambridge University Press.

\bibitem[Weinstein(1996)]{Weinstein1996}
Weinstein, A.\ (1996).
Groupoids: Unifying internal and external symmetry.
\textit{Notices of the AMS} \textbf{43}(7):744--752.

\bibitem[AWMLG(2026)]{AwmlgLiverpool2026}
Anonymous.\ (2026).
Adaptive world model learning for portfolio optimization.
\textit{Preprint}.

\bibitem[JPMorgan(2019)]{MorganJPMaGAN}
Cont, R., Stoikov, S., and Talreja, R.\ (2010).
A stochastic model for order book dynamics.
\textit{Operations Research} \textbf{58}(3):549--563.

\end{thebibliography}

\end{document}
